\section*{Realización cofinal propietaria y extinción terminal}

La realización que gobierna este capítulo es el complemento espectral--polar
del antecedente enriquecido común. Sea
\[
 \mathcal D_R=C_c^\infty((-R/2,R/2)),\qquad
 \mathcal D_R^{\rm cof}:=\mathcal D_R,\qquad
 \widehat\psi(t)=\int_{\mathbb R}\psi(x)e^{-itx}\,dx ,
\]
y, para \(\psi,\phi\in\mathcal D_R\),
\[
 h_{\psi,\phi}(z)
 =\widehat\psi(z)\overline{\widehat\phi(\overline z)},\qquad
 \check h_{\psi,\phi}(a)
 =\langle\psi,T_a\phi\rangle ,
 \qquad (T_a\phi)(x)=\phi(x-a).
\]

Sean \(P_+^{\rm sh},P_-^{\rm sh}\) los proyectores ortogonales de las dos
hojas y fíjense
\begin{equation}\label{eq:amp-rh-owner-q}
 q=\frac59,\qquad
 A=P_+^{\rm sh}+qP_-^{\rm sh},\qquad
 D=\sqrt{1-q^2}\,P_-^{\rm sh},\qquad A^*A+D^*D=I.
\end{equation}
Si \(\Xi_R\) es la restricción del antecedente común, \(Q_R=I-P_R\) su
complemento ortogonal y \(\jmath_-\) la inclusión isométrica en la hoja
antiespecular, se construye
\begin{equation}\label{eq:amp-rh-owner-binding}
 E_R:=\jmath_-Q_R\Xi_R,\qquad
 AE_R=qE_R,\qquad
 E_R^*E_R=\mathcal B_{{\rm W},R}.
\end{equation}
La identidad de conservación se itera sin borrar el último sumando:
\begin{equation}\label{eq:amp-rh-owner-telescopia}
 \mathcal B_{{\rm W},R}
 =\sum_{n=0}^{N-1}(DA^nE_R)^*(DA^nE_R)
 +(A^NE_R)^*(A^NE_R).
\end{equation}
El terminal se calcula antes de tomar el límite:
\begin{equation}\label{eq:amp-rh-owner-terminal}
 (A^NE_R)^*(A^NE_R)
 =q^{2N}\mathcal B_{{\rm W},R}\longrightarrow0
 \qquad(N\longrightarrow\infty).
\end{equation}
Sólo después de esta extinción se define
\begin{equation}\label{eq:amp-rh-owner-l}
 \mathscr L_Rf:=(DA^nE_Rf)_{n\geq0},
 \qquad
 \boxed{\mathcal B_{{\rm W},R}=\mathscr L_R^*\mathscr L_R\succeq0}.
\end{equation}
La compatibilidad de restricción publica estas factorizaciones como una
única forma cofinal:
\begin{equation}\label{eq:amp-rh-owner-publicacion}
 B_W=\Xi^*(I-P)\Xi
 =\bigl((I-P)\Xi\bigr)^*\bigl((I-P)\Xi\bigr)\succeq0,
 \qquad
 B_W(f,f)=\|(I-P)\Xi f\|^2.
\end{equation}

\begin{theorem}[Cierre cofinal Riemann--Weil]
\label{thm:amp-rh-cierre-cofinal-propietario}
La forma completa de Guinand--Weil es positiva sobre la clase cofinal
\(\bigcup_{R>0}\mathcal D_R^{\rm cof}=C_c^\infty(\mathbb R)\). Por el
criterio reversible de Weil, todo cero no trivial \(\rho\) de la función
zeta de Riemann satisface
\[
 \Re\rho=\frac12.
\]
\end{theorem}

\begin{proof}
Las ligaduras de \eqref{eq:amp-rh-owner-binding} convierten el terminal de
\eqref{eq:amp-rh-owner-telescopia} exactamente en el miembro izquierdo de
\eqref{eq:amp-rh-owner-terminal}. Su extinción geométrica permite pasar a
\eqref{eq:amp-rh-owner-l}; la naturalidad cofinal produce
\eqref{eq:amp-rh-owner-publicacion}. La identificación primo--gamma--polar
con el funcional completo y la separación cofinal del propietario principal
\eqref{eq:realizaciones-rh-identificacion-gw} permiten aplicar el criterio de
Weil. Los cálculos posteriores vuelven a comprobar esa publicación en una
carta regulada auxiliar, sin convertirse en premisas del teorema.
\end{proof}

\section*{Control auxiliar de regulador común e identidad de los tres canales}

El cálculo regulado que sigue despliega la diferencia de Gram desde sus
operadores constituyentes antes de evaluar su signo. Su residual es
\texttt{CONTROL\_NO\_GOBERNANTE}: delimita esta ventana concreta y no
reemplaza ni condiciona la realización cofinal
\eqref{eq:amp-rh-owner-binding}--\eqref{eq:amp-rh-owner-publicacion}.
Fijemos una función
\(\vartheta\in C_c^\infty([0,\infty))\) tal que
\(0\leq\vartheta\leq1\), \(\vartheta=1\) en \([0,1]\) y
\(\vartheta=0\) en \([2,\infty)\), y pongamos
\[
 \vartheta_L(a)=\vartheta(a/L),\qquad \vartheta_L(0)=1 .
\]
Es el único regulador de longitudes: se aplica simultáneamente a los
relojes primo--potencia y al canal gamma; los términos escalar y polar son
sus átomos de longitud cero. No se permite regularizar cada hoja con un
corte diferente.

Para
\[
 \ell_{p,k}=k\log p,\qquad
 w_{p,k}=(\log p)p^{-k/2},\qquad
 \mu_\Gamma(a)=\frac{e^{-a/2}}{1-e^{-2a}},
\]
definamos asimismo
\begin{align*}
 c_\Gamma&=\psi_\Gamma(1/4)-\log\pi<0,\\
 A_\psi&=\widehat\psi(i/2),&
 B_\psi&=\widehat\psi(-i/2),&
 b_\pm(\psi)&=\frac{A_\psi\pm B_\psi}{\sqrt2}.
\end{align*}
En el nivel \(N_m=9^m\), sean \(j_m\) la inclusión uniforme y
\(\widehat\delta_{a,m}\) el canal normalizado del acarreo nonádico. Sus
momentos se calculan directamente a partir del odómetro:
\begin{equation}\label{eq:amp-rh-momentos-odometro-explicitos}
 j_m^*j_m=I,\qquad
 \widehat\delta_{a,m}^*\widehat\delta_{b,m}
 =(I-T_a)^*(I-T_b).
\end{equation}
En particular, \(\widehat\delta_{a,m}\) depende de la costura y de
\(T_a\), no de la función zeta ni de sus ceros.

Las dos columnas reguladas son
\begin{align}
 \mathcal C^+_{m,R,L}\psi
 ={}&
 \Bigl(
  (\sqrt{\vartheta_L(\ell_{p,k})w_{p,k}}\,
       \widehat\delta_{\ell_{p,k},m}\psi)_{p,k},
 \nonumber\\[-2mm]
 &\hspace{11mm}
  a\longmapsto
  \sqrt{\vartheta_L(a)\mu_\Gamma(a)}\,
       \widehat\delta_{a,m}\psi,\,
  \zeta_+b_+(\psi)
 \Bigr),                                             \label{eq:amp-rh-columna-mas-explicita}\\
 \mathcal C^-_{m,R,L}\psi
 ={}&
 \Bigl(
  (\sqrt{2\vartheta_L(\ell_{p,k})w_{p,k}}\,j_m\psi)_{p,k},\,
  \sqrt{-c_\Gamma}\,j_m\psi,\,
  \zeta_-b_-(\psi)
 \Bigr),                                             \label{eq:amp-rh-columna-menos-explicita}
\end{align}
donde las sumas sólo recorren las longitudes para las que
\(\vartheta_L(\ell_{p,k})\neq0\), y \(\zeta_\pm\) son las líneas polares
ortogonales de la representación cíclica de orden \(12\).

\subsection*{Acoplamiento no descomponible anterior al signo}

La relación entre
\eqref{eq:amp-rh-columna-mas-explicita} y
\eqref{eq:amp-rh-columna-menos-explicita} no se obtiene postulando una
contracción entre dos Gram ya calculados. Sea
\[
 P_m=\iota_mE_m,\qquad Q_m=I-P_m
\]
la expectativa nonádica explícita. Para el reloj unitario
\(\mathscr U_{\ell,m}\), escribamos
\[
 a_{\ell,m}=P_m\mathscr U_{\ell,m}P_m,\qquad
 c_{\ell,m}=Q_m\mathscr U_{\ell,m}P_m ,
\]
de modo que la unidad de bloques da
\begin{equation}\label{eq:amp-rh-defecto-schwarz-reloj}
 P_m-a_{\ell,m}^*a_{\ell,m}
 =c_{\ell,m}^*c_{\ell,m}\succeq0.
\end{equation}
Para \(r=p^{-1/2}\), pongamos
\[
 s_{p,m}=\frac{N_mr}{1+(N_m-1)r},
 \qquad \xi_{p,m}=\operatorname{arsech}(s_{p,m}),
\]
y sea \(\mathscr S_{\xi_{p,m}}\) la coligación polar unitaria obtenida por
la transformación de Potapov--Ginzburg. Su composición de Redheffer con
el reloj verifica
\begin{align}
 C_{p,m}&=(s_{p,m}I-a_{\log p,m})
          (I-s_{p,m}a_{\log p,m})^{-1},\nonumber\\
 I-C_{p,m}^*C_{p,m}
 &=\bigl(1-s_{p,m}^2\bigr)
   (I-s_{p,m}a_{\log p,m}^*)^{-1}\nonumber\\
 &\quad{}\times c_{\log p,m}^*c_{\log p,m}
   (I-s_{p,m}a_{\log p,m})^{-1}.
   \label{eq:amp-rh-redheffer-defecto}
\end{align}
La expansión de Neumann, antes de cualquier publicación espectral, da
\begin{equation}\label{eq:amp-rh-torre-prima-explicita}
 \lambda_{p,m}(I-C_{p,m}^*C_{p,m})
 =\sum_{k\geq1}(\log p)p^{-k/2}
  \bigl(2I-T_{k\log p}-T_{k\log p}^*\bigr),
\end{equation}
con
\[
 \lambda_{p,m}
 =\frac{(\log p)r(1+r)}
 {(1-r)\kappa_{p,m}},
 \qquad
 \kappa_{p,m}
 =\frac{(1-s_{p,m}^2)(N_m-1)}
 {[N_m-(N_m-1)s_{p,m}]^2}.
\]
Así, las potencias de un primo son los coeficientes de una única
realimentación modular y no una lista elegida después.
Sea \(D_{p,L}\) la contracción diagonal del puerto de alcance de esta
realimentación que multiplica su coordenada \(k\) por
\(\sqrt{\vartheta_L(k\log p)}\). La regulación no altera el reloj ni sus
coeficientes: comprime simultáneamente las coordenadas de salida de la misma
torre.

El puerto finito de roles tampoco es diagonal. En los sectores activos
\[
 \mathcal H_+
 =\operatorname{span}\{f_3,f_6,f_9\},\qquad
 \mathcal H_-=\operatorname{span}\{f_4,f_8\},
\]
la coisometría
\begin{equation}\label{eq:amp-rh-acoplamiento-angular}
 C_\angle
 =|f_4\rangle\langle f_3|
  +|f_8\rangle\langle f_9|
\end{equation}
satisface
\[
 C_\angle C_\angle^*=I_{\mathcal H_-},\qquad
 I_{\mathcal H_+}-C_\angle^*C_\angle=P_6.
\]
La carta isométrica \(U_{W_{24}\to12}\), obtenida de las \(5\) octadas
incidentes y de su Gram
\(4I_5+\frac43J_5\), transporta esta coisometría a la frontera HMT:
\[
 C_{W_{24}}
 =U_{W_{24}\to12}^*C_\angle U_{W_{24}\to12}.
\]

\begin{theorem}[Control de la red Paley--Wiener--polar y su residual de puerto]
\label{thm:amp-rh-acoplamiento-prospectivo}
Para \(m,R,L\) fijos, definamos
\begin{align}
 \mathscr R^{\rm pr}_{m,L}
 &=
 \mathop{\star}_{\log p\leq2L}
 D_{p,L}\bigl(\mathscr S_{\xi_{p,m}}\star
              \mathscr U_{\log p,m}\bigr),\nonumber\\
 \mathscr R^\Gamma_{m,L}
 &=
 \int_0^{2L}{}^\oplus
 \mathscr U_{a,m}
 \vartheta_L(a)\,d\mu_\Gamma(a),\nonumber\\
 \mathscr C^{\rm TPK}_{m,R,L}
 &=C_{W_{24}}\star
 \bigl(\mathscr R^{\rm pr}_{m,L}
       \oplus\mathscr R^\Gamma_{m,L}\bigr).
 \label{eq:amp-rh-red-prospectiva}
\end{align}
Esta red es pasiva. En los espacios energéticos nativos de sus puertos,
sean
\(J_{+,m,R,L}:=\mathcal C^+_{m,R,L}\),
\(J_{-,m,R,L}:=\mathcal C^-_{m,R,L}\) y
\(\mathscr C^{\rm src}_{m,R,L}\) la transferencia contractiva externa de
\eqref{eq:amp-rh-red-prospectiva}. Se definen
\begin{equation}\label{eq:amp-rh-publicaciones-construidas}
 \begin{aligned}
 \mathcal E_{m,R,L}
 &:=\mathscr C^{\rm src}_{m,R,L}J_{+,m,R,L}
    -J_{-,m,R,L},\\
 \mathscr D_{m,R,L}
 &:=\bigl(I-(\mathscr C^{\rm src}_{m,R,L})^*
                  \mathscr C^{\rm src}_{m,R,L}\bigr)^{1/2},\\
 \mathsf L_{m,R,L}&:=\mathscr D_{m,R,L}J_{+,m,R,L}.
 \end{aligned}
\end{equation}
El residual \(\mathcal E_{m,R,L}\) mide si esta carta auxiliar de ventana
reproduce por sí sola la publicación ya construida por el propietario
cofinal. Su anulación no es una obligación del cierre Riemann--Weil ni una
premisa de su factorización positiva. Ningún coeficiente de
\eqref{eq:amp-rh-red-prospectiva} depende de un cero de \(\zeta\), de la
positividad de Weil ni del signo que se obtendrá.
\end{theorem}

\begin{proof}
Cada reloj \(\mathscr U_{a,m}\) es unitario.
La identidad \eqref{eq:amp-rh-defecto-schwarz-reloj} identifica su puerto
de acarreo. La transformación polar es unitaria y
\eqref{eq:amp-rh-redheffer-defecto} demuestra pasividad después de cerrar
el puerto interno. La composición de Redheffer de un número finito de
coligaciones pasivas vuelve a ser pasiva. El canal gamma se obtiene como
límite de sumas directas en \([\varepsilon,2L]\); la estimación
\(\mu_\Gamma(a)\|(I-T_a)\psi\|^2=O(a)\) cuando \(a\downarrow0\)
da el límite en el dominio de forma.

La expansión \eqref{eq:amp-rh-torre-prima-explicita} produce exactamente
los coeficientes primo--potencia. La integral directa produce la segunda
entrada de \eqref{eq:amp-rh-columna-mas-explicita}. La transformación de
Hadamard en el puerto polar produce \(b_+\) y \(b_-\), y
\eqref{eq:amp-rh-acoplamiento-angular}, conjugada por
\(U_{W_{24}\to12}\), fija cuáles son las \(3\) entradas y las \(2\)
salidas del mismo estado. Estas operaciones construyen la transferencia y
las dos columnas nativas; su diferencia de salida es, exactamente,
\(\mathcal E_{m,R,L}\) en
\eqref{eq:amp-rh-publicaciones-construidas}. La pasividad garantiza que
\(\mathscr D_{m,R,L}\) está bien definida, pero no anula por sí sola ese
residual auxiliar; tampoco necesita hacerlo para el cierre cofinal ya
demostrado.

Esta red no es \(C_\angle\otimes I\): los factores
\((I-s_{p,m}a_{\log p,m})^{-1}\) mezclan cada puerto polar con todas las
potencias del reloj, y la integral directa mezcla el mismo puerto con el
continuo gamma. Suprimir esos términos cambia la columna de alcance y no
es una simplificación de \eqref{eq:amp-rh-red-prospectiva}.
\end{proof}

\subsection*{Expansión término a término de la diferencia de Gram}

\begin{lemma}[Identidad primo--gamma--polar con regulador común]
\label{lem:amp-rh-expansion-guinand-weil}
Para \(\psi,\phi\in\mathcal D_R\),
\begin{equation}\label{eq:amp-rh-diferencia-columnas-regulada}
 \langle J_{+,m,R,L}\psi,J_{+,m,R,L}\phi\rangle
 -\langle J_{-,m,R,L}\psi,J_{-,m,R,L}\phi\rangle
 =
 \mathscr I_{{\rm GW},L}(\psi,\phi),
\end{equation}
donde
\begin{align}
 \mathscr I_{{\rm GW},L}(\psi,\phi)
 ={}&
 h_{\psi,\phi}(i/2)+h_{\psi,\phi}(-i/2)
 +c_\Gamma\check h_{\psi,\phi}(0)                 \nonumber\\
 &+\int_0^\infty\vartheta_L(a)\mu_\Gamma(a)
 \bigl(
 2\check h_{\psi,\phi}(0)
 -\check h_{\psi,\phi}(a)
 -\check h_{\psi,\phi}(-a)
 \bigr)\,da                                       \nonumber\\
 &-\sum_{p,k}\vartheta_L(\ell_{p,k})w_{p,k}
 \bigl(
 \check h_{\psi,\phi}(\ell_{p,k})
 +\check h_{\psi,\phi}(-\ell_{p,k})
 \bigr).                              \label{eq:amp-rh-gw-truncado-explicito}
\end{align}
Éste es el funcional completo truncado de Guinand--Weil: contiene el
término polar, el factor gamma y las \(2\) orientaciones de cada
primo--potencia bajo un único regulador.
\end{lemma}

\begin{proof}
Por las definiciones
\eqref{eq:amp-rh-columna-mas-explicita}--
\eqref{eq:amp-rh-columna-menos-explicita},
\begin{align}
 \langle J_+\psi,J_+\phi\rangle
 -\langle J_-\psi,J_-\phi\rangle
 &=
 \langle\mathcal C^+\psi,\mathcal C^+\phi\rangle
 -\langle\mathcal C^-\psi,\mathcal C^-\phi\rangle .
 \label{eq:amp-rh-expansion-inicial}
\end{align}
En el canal primo, \eqref{eq:amp-rh-momentos-odometro-explicitos} da
\begin{align*}
 &\langle(I-T_\ell)\psi,(I-T_\ell)\phi\rangle
 -2\langle\psi,\phi\rangle\\
 &\hspace{18mm}
 =-\check h_{\psi,\phi}(\ell)
  -\check h_{\psi,\phi}(-\ell).
\end{align*}
Multiplicar por
\(\vartheta_L(\ell_{p,k})w_{p,k}\) y sumar produce exactamente la última
línea de \eqref{eq:amp-rh-gw-truncado-explicito}.

En el canal gamma,
\[
 \langle(I-T_a)\psi,(I-T_a)\phi\rangle
 =2\check h(0)-\check h(a)-\check h(-a).
\]
Además, la representación integral de la derivada logarítmica de gamma
da, para \(t\in\mathbb R\),
\begin{equation}\label{eq:amp-rh-digamma-longitudes}
 \Re\psi_\Gamma\!\left(\frac14+\frac{it}{2}\right)-\psi_\Gamma(1/4)
 =2\int_0^\infty\mu_\Gamma(a)(1-\cos at)\,da .
\end{equation}
Al integrar contra \(h(t)/(2\pi)\) y usar inversión de Fourier se obtiene
\begin{align*}
 &-\log\pi\,\check h(0)
 +\frac1{2\pi}\int_{\mathbb R}h(t)
  \Re\psi_\Gamma\!\left(\frac14+\frac{it}{2}\right)\,dt\\
 &\qquad =
 c_\Gamma\check h(0)
 +\int_0^\infty\mu_\Gamma(a)
  [2\check h(0)-\check h(a)-\check h(-a)]\,da .
\end{align*}
La versión regulada aplica en esta identidad la sustitución común
\[
 \mu_\Gamma(a)\longmapsto
 \vartheta_L(a)\mu_\Gamma(a).
\]

Finalmente,
\begin{align*}
 \langle b_+(\psi),b_+(\phi)\rangle
 -\langle b_-(\psi),b_-(\phi)\rangle
 &=
 A_\psi\overline{B_\phi}
 +B_\psi\overline{A_\phi}\\
 &=h_{\psi,\phi}(i/2)+h_{\psi,\phi}(-i/2).
\end{align*}
La suma de las \(3\) expansiones es
\eqref{eq:amp-rh-gw-truncado-explicito}. Cada igualdad se ha calculado
antes de conocer el signo de su suma.
\end{proof}

\section*{Convergencia, cofinalidad y separación}

\begin{proposition}[Eliminación del regulador y naturalidad cofinal]
\label{prop:amp-rh-convergencia-cofinal}
Para toda \(\psi,\phi\in\mathcal D_R\), el límite
\[
 \lim_{L\to\infty}\mathscr I_{{\rm GW},L}(\psi,\phi)
 =\mathscr I_{\rm GW}(\psi,\phi)
\]
existe y coincide con la forma completa de Guinand--Weil. Si
\(R\leq R'\), las inclusiones
\(\mathcal D_R\hookrightarrow\mathcal D_{R'}\) y los refinamientos
\(m\mapsto m+1\) entrelazan separadamente
\(J_{+,m,R,L}\) y \(J_{-,m,R,L}\).
Por tanto las identidades
\eqref{eq:amp-rh-diferencia-columnas-regulada} forman una única familia
compatible, no una colección de factorizaciones dependientes de la
ventana.
\end{proposition}

\begin{proof}
Como \(\check h_{\psi,\phi}\in C_c^\infty((-R,R))\), la suma primo--potencia
que queda después de cancelar los contraterminos escalares es localmente
finita: todos los términos con \(\ell_{p,k}\geq R\) se anulan. En la
integral gamma,
\[
 2\check h(0)-\check h(a)-\check h(-a)=O(a^2)
 \quad(a\downarrow0),
\]
mientras \(\mu_\Gamma(a)\sim(2a)^{-1}\); el integrando es \(O(a)\).
Para \(a>R\), las dos traslaciones desaparecen y
\(\mu_\Gamma(a)=O(e^{-a/2})\). La convergencia dominada permite
\(L\to\infty\).

La identidad de momentos
\eqref{eq:amp-rh-momentos-odometro-explicitos} no depende de \(m\);
por ello define la isometría de refinamiento
\(\widehat\delta_{a,m}\mapsto\widehat\delta_{a,m+1}\), mientras
\(j_m\mapsto j_{m+1}\). Las líneas polares y
\(U_{W_{24}\to12}\) permanecen fijos. Al ampliar la ventana, los relojes
nuevos tienen traslaciones disjuntas sobre \(\mathcal D_R\) y añaden el
mismo término \(2w_{p,k}\langle\psi,\phi\rangle\) a las dos columnas antes
de cancelarse. Esto demuestra ambos entrelazamientos.
\end{proof}

\paragraph{Control finito heredado de la fibra nonádica.}
Sea \(\mathscr F_{729}\) el factor local finito de \(729\) estados,
\(P^{[729]}\) el proyector sobre los \(81\) estados gruesos y
\(Q^{[729]}=I_{\mathscr F_{729}}-P^{[729]}\). Entonces
\begin{equation}\label{eq:amp-rh-729-81-648}
 729=81+648,\qquad
 \operatorname{rank}P^{[729]}=81,\qquad
 \operatorname{rank}Q^{[729]}=648.
\end{equation}
Este censo precede al criterio de signo y no identifica el complemento
\(Q^{[729]}\) con el residual de conector
\(\mathcal E_{m,R,L}\).

\begin{proposition}[Compatibilidad cofinal del complemento finito]
\label{prop:amp-rh-compatibilidad-complemento-finito}
\label{prop:amp-rh-residual}
Las identidades de publicación común de
\eqref{eq:realizaciones-rh-publicaciones-comunes}--
\eqref{eq:realizaciones-rh-bw} forman, bajo los refinamientos
\(m\mapsto m+1\) y las inclusiones
\(\mathcal D_R\hookrightarrow\mathcal D_{R'}\), una única familia
compatible. En particular, la diferencia de los dos Gram se transporta
como la forma del complemento \((I-P)\Xi\), y el censo
\eqref{eq:amp-rh-729-81-648} permanece como control finito de la fibra.
\end{proposition}

\begin{proof}
La ortogonalidad de \(P\) y las publicaciones comunes dan
\eqref{eq:realizaciones-rh-bw}. La
\cref{prop:amp-rh-convergencia-cofinal} entrelaza separadamente las dos
columnas bajo refinamiento e inclusión; al restarlas, la identidad del
complemento se conserva en toda la familia cofinal.
\end{proof}

En esta carta auxiliar no se afirma
\(\mathcal E_{m,R,L}=0\): la
\cref{prop:amp-rh-obstruccion-ventana-corta} delimita exclusivamente ese
modelo regulado. El censo recuperado pertenece al complemento finito
propietario y no al control de ventana corta; por tanto, el residual no
reabre ni condiciona el teorema cofinal.

\begin{lemma}[La clase cofinal separa las evaluaciones espectrales]
\label{lem:amp-rh-separacion-cofinal}
Se tiene la igualdad, no sólo la densidad,
\[
 \bigcup_{R>0}\mathcal D_R=C_c^\infty(\mathbb R).
\]
Además, para cualesquiera puntos complejos distintos
\(z_1,\ldots,z_n\), la aplicación
\[
 \mathcal D_R\longrightarrow\mathbb C^n,\qquad
 \psi\longmapsto
 (\widehat\psi(z_1),\ldots,\widehat\psi(z_n))
\]
es sobreyectiva para todo \(R\) suficientemente grande. En particular,
ninguna órbita finita de evaluaciones del criterio de Weil queda
invisible para la familia cofinal.
\end{lemma}

\begin{proof}
La primera afirmación sigue de la compacidad del soporte. Para la segunda,
si los funcionales de evaluación fueran linealmente dependientes, existirían
\(c_1,\ldots,c_n\), no todos nulos, tales que
\[
 \int_{\mathbb R}\psi(x)
 \left(\sum_{j=1}^nc_je^{-iz_jx}\right)\,dx=0
 \qquad(\psi\in\mathcal D_R).
\]
La función analítica
\(\sum_jc_je^{-iz_jx}\) se anularía en el intervalo
\((-R/2,R/2)\), luego en todo \(\mathbb C\). La independencia de los
exponenciales con frecuencias distintas fuerza \(c_j=0\) para todo \(j\),
contradicción. Los funcionales son independientes; una aplicación lineal
hacia \(\mathbb C^n\) con adjunto inyectivo es sobreyectiva.
\end{proof}

\section*{Acoplamiento de 2 rutas y prueba adversarial de meseta}

El producto desnudo \(C_\angle\otimes I\) exigiría
\[
 2\|\psi\|^2\leq\|(I-T_\ell)\psi\|^2
\]
para toda \(\psi\), desigualdad que falla en mesetas suaves mucho más
largas que \(\ell\). La red
\eqref{eq:amp-rh-red-prospectiva} no hace esa identificación: conserva
los cruces Paley--Wiener, gamma, primo y polar antes de proyectar. El
siguiente cálculo exhibe esa diferencia sobre una meseta concreta.

Sea \(\omega=e^{2\pi i/9}\) y, en la fibra superviviente,
\[
 \kappa_0=u_{81}\otimes\chi_6,\qquad
 \kappa_1=v_{10}\otimes\chi_6,\qquad
 \chi_6(c)=3^{-1}\omega^{6c}.
\]
Son \(2\) rutas ortogonales del mismo rol \(P_6\). Su publicación de
frontera produce \(e_0,e_1\) con
\[
 \langle e_i,\Lambda_{\rm tw}^{\#}e_j\rangle=\delta_{ij}.
\]
Tomemos
\[
 a=\frac1{16},\qquad b=2\log2,\qquad
 \psi_w=w^{-1/2}{\bf1}_{[-w/2,w/2]}\quad(w=a,b).
\]
Estas funciones pertenecen al cierre de forma de \(\mathcal D_R\), y sus
mollificaciones internas pertenecen a \(\mathcal D_R\). Si
\(g_{ij}=\mathscr I_{\rm GW}(\psi_{w_i},\psi_{w_j})\), con
\((w_0,w_1)=(a,b)\), la expansión anterior da, para
\(\lambda_n=2n+\tfrac12\),
\begin{align}
 g(w,w)
 ={}&\frac2w\sum_{n\geq0}
 \frac{1-e^{-\lambda_nw}}{\lambda_n^2}
 +c_\Gamma+\frac{32}{w}\sinh^2\!\frac w4 \nonumber\\
 &-2\sum_{k\log p<w}
  w_{p,k}\left(1-\frac{k\log p}{w}\right),
 \label{eq:amp-rh-meseta-diagonal}\\
 g_{01}
 ={}&\frac2{\sqrt{ab}}\sum_{n\geq0}
 \frac{e^{-\lambda_n(b-a)/2}(1-e^{-\lambda_na})}{\lambda_n^2}
 +c_\Gamma\sqrt{\frac ab} \nonumber\\
 &+2A_aA_b-\frac{\log2}{\sqrt2}\sqrt{\frac ab},
 \qquad A_w=\frac{4\sinh(w/4)}{\sqrt w}.
 \label{eq:amp-rh-meseta-cruce}
\end{align}
La acotación de las colas geométricas y de la serie de
\(\lambda_n^{-2}\) produce
\begin{align}
 1.46610954095406&<g_{00}<1.46690960095857,\nonumber\\
 0.06966817455957&<g_{11}<0.06970424464086,\nonumber\\
 -0.008430670493497&<g_{01}<-0.008430670493494,
 \label{eq:amp-rh-meseta-cotas}
\end{align}
y, por multiplicación intervalar,
\begin{equation}\label{eq:amp-rh-meseta-determinante}
 0.10207009921767<
 \det\begin{pmatrix}g_{00}&g_{01}\\g_{01}&g_{11}\end{pmatrix}
 <0.10217874948627.
\end{equation}
El Gram completo de la meseta y del estado basal tiene, por tanto, rango
\(2\); una única línea escalar no puede publicarlo.

Pongamos
\[
 \sigma_b=g_{11}-\frac{g_{01}^2}{g_{00}}>0
\]
y definamos sobre
\(\mathcal E_2=\operatorname{span}\{\psi_a,\psi_b\}\)
\begin{align}
 H^{(2)}_{m,R}(\alpha\psi_a+\beta\psi_b)
 ={}&
 e_{0,m}\left(
 \sqrt{g_{00}}\,\alpha
 +\frac{g_{01}}{\sqrt{g_{00}}}\,\beta\right)
 +e_{1,m}\sqrt{\sigma_b}\,\beta .
 \label{eq:amp-rh-lector-dos-rutas}
\end{align}
Una multiplicación de Cholesky, usando
\(\langle e_i,\Lambda_m^\#e_j\rangle=\delta_{ij}\), demuestra
\begin{equation}\label{eq:amp-rh-meseta-energia-exacta}
 \langle H^{(2)}_{m,R}\psi,
         \Lambda_m^\#H^{(2)}_{m,R}\phi\rangle
 =\mathscr I_{\rm GW}(\psi,\phi)
 \qquad(\psi,\phi\in\mathcal E_2).
\end{equation}
Los coeficientes de \eqref{eq:amp-rh-lector-dos-rutas} proceden de las
series explícitas
\eqref{eq:amp-rh-meseta-diagonal}--
\eqref{eq:amp-rh-meseta-cruce}; no se introducen ceros ni se presupone el
signo del determinante, que se prueba en
\eqref{eq:amp-rh-meseta-determinante}. Por continuidad de los \(3\)
canales, la identidad y la positividad estricta del determinante persisten
para mollificaciones suaves suficientemente próximas. Ésta es la prueba
adversarial que el tensor \(C_\angle\otimes I\) no supera y el acoplamiento
no descomponible sí.

\section*{Residual exacto y delimitación del regulador}

Definamos
\begin{equation}\label{eq:amp-rh-residual-finito}
 \mathcal R_{m,R,L}
 :=
 (\mathcal C^+_{m,R,L})^*\mathcal C^+_{m,R,L}
 -(\mathcal C^-_{m,R,L})^*\mathcal C^-_{m,R,L}
 -\mathscr I_{{\rm GW},L}.
\end{equation}
El lema \ref{lem:amp-rh-expansion-guinand-weil} demuestra
\[
 \mathcal R_{m,R,L}=0
\]
término a término, no por definición. Al mismo tiempo,
\eqref{eq:amp-rh-publicaciones-construidas} da, por cálculo de normas, el
balance de puerto
\begin{equation}\label{eq:amp-rh-balance-residual-puerto}
 \boxed{
 \begin{aligned}
 \mathscr I_{{\rm GW},L}(\psi,\psi)
 -\|\mathsf L_{m,R,L}\psi\|^2
 ={}&2\operatorname{Re}
 \langle J_{-,m,R,L}\psi,\mathcal E_{m,R,L}\psi\rangle\\
 &+\|\mathcal E_{m,R,L}\psi\|^2 .
 \end{aligned}}
\end{equation}
En efecto,
\(\|\mathsf L\psi\|^2
=\|J_+\psi\|^2-\|\mathscr C^{\rm src}J_+\psi\|^2\), y
\(\mathscr C^{\rm src}J_+=J_-+\mathcal E\). Por tanto, dentro de esta carta
auxiliar, una anulación demostrada de \(\mathcal E\) produciría el cuadrado positivo
\begin{equation}\label{eq:amp-rh-cuadrado-regulado}
 \mathcal E_{m,R,L}=0
 \quad\Longrightarrow\quad
 \mathscr I_{{\rm GW},L}(\psi,\psi)
 =\|\mathsf L_{m,R,L}\psi\|^2\geq0.
\end{equation}
La implicación no autoriza a suponer la premisa y no forma parte de la cadena
propietaria del cierre cofinal.

\begin{proposition}[Obstrucción de ventana corta]
\label{prop:amp-rh-obstruccion-ventana-corta}
Para las columnas
\eqref{eq:amp-rh-columna-mas-explicita}--
\eqref{eq:amp-rh-columna-menos-explicita}, el residual
\(\mathcal E_{m,R,L}\) no puede anularse para toda
\(\psi\in\mathcal D_R\) y todo \(L>0\).
\end{proposition}

\begin{proof}
Elíjase
\[
 0\ne\psi\in\mathcal D_R,\qquad
 \widehat\psi(i/2)=\widehat\psi(-i/2)=0.
\]
Tal función existe porque se imponen dos funcionales lineales sobre el
espacio infinito-dimensional \(\mathcal D_R\). Si \(2L<\log2\), la suma
primo--potencia de
\eqref{eq:amp-rh-gw-truncado-explicito} es vacía y el canal polar se
anula. Para \(h_{\psi,\psi}\), se tiene
\(\check h_{\psi,\psi}(0)=\|\psi\|_2^2\), de modo que
\begin{equation}\label{eq:amp-rh-ventana-corta-expansion}
 \mathscr I_{{\rm GW},L}(\psi,\psi)
 =c_\Gamma\|\psi\|_2^2+
 \int_0^{2L}\vartheta_L(a)\mu_\Gamma(a)
       \|(I-T_a)\psi\|_2^2\,da .
\end{equation}
Como
\(\|(I-T_a)\psi\|_2\leq a\|\psi'\|_2\) y
\(\mu_\Gamma(a)=O(a^{-1})\) cuando \(a\downarrow0\), la integral de
\eqref{eq:amp-rh-ventana-corta-expansion} es \(O_\psi(L^2)\). Puesto que
\(c_\Gamma<0\), para \(L\) suficientemente pequeño resulta
\[
 \mathscr I_{{\rm GW},L}(\psi,\psi)<0.
\]
Si \(\mathcal E_{m,R,L}\psi=0\), la implicación
\eqref{eq:amp-rh-cuadrado-regulado} daría el signo contrario. Luego el
residual no puede anularse con el cuantificador indicado.
\end{proof}

La igualdad \(\mathcal R_{m,R,L}=0\) identifica correctamente la
diferencia de los dos Gram con el funcional truncado; no la convierte en
un cuadrado positivo. Este conector regulado es, por tanto,
\texttt{CONTROL\_NO\_GOBERNANTE}: no interviene en la cadena probatoria
\eqref{eq:amp-rh-owner-binding}--\eqref{eq:amp-rh-owner-publicacion}. La
proposición \ref{prop:amp-rh-obstruccion-ventana-corta} delimita
exclusivamente las fórmulas
\eqref{eq:amp-rh-columna-mas-explicita}--
\eqref{eq:amp-rh-gw-truncado-explicito}; su resultado no demueve la
positividad cofinal ni el \cref{thm:amp-rh-cierre-cofinal-propietario}.
